\documentclass[10pt]{article}
\usepackage[T1]{fontenc}
\usepackage[utf8]{inputenc}
\usepackage[french]{babel}
\usepackage{amsmath,amssymb,mathrsfs,mathpazo}
\usepackage{geometry}
\geometry{a4paper,margin=2cm}
\usepackage{graphicx,array,tabularx,tabularray,colortbl,multirow,multicol,booktabs,enumitem,pifont,ulem,fancybox,diagbox,cancel}
\usepackage{pstricks,pstricks-add,pst-plot,pst-tree,pst-node,pst-text,pst-func,pst-math,pst-bspline,pst-3dplot,pst-eucl,pst-solides3d}
\let\pscalculate\undefined
\usepackage{pst-bezier}
\usepackage{tikz,pgfplots,tkz-tab}
\usetikzlibrary{arrows,arrows.meta,patterns,calc,positioning,trees,decorations.markings,shapes.geometric,shapes.misc,shapes.arrows,intersections,decorations.pathmorphing,decorations.pathreplacing,backgrounds,shadows,babel}
\pgfplotsset{compat=1.18}
\usepgfplotslibrary{fillbetween}
\usepackage[np]{numprint}
\usepackage{listings,makecell,pas-tableur,twoopt}
\let\barre\undefined
\usepackage[pstricks]{tablvar}
\definecolor{darkgreen}{RGB}{0,100,0}
\providecommand{\red}{\color{red}}
\providecommand{\blue}{\color{blue}}
\providecommand{\green}{\color{green}}
\providecommand{\black}{\color{black}}
\providecommand{\white}{\color{white}}
\providecommand{\gray}{\color{gray}}
\providecommand{\cyan}{\color{cyan}}
\providecommand{\magenta}{\color{magenta}}
\providecommand{\yellow}{\color{yellow}}
\providecommand{\violet}{\color{violet}}
\providecommand{\orange}{\color{orange}}
\providecommand{\euro}{\texteuro}
\providecommand{\EUR}{\texteuro}
\providecommand{\ieme}{\textsuperscript{e}}
\providecommand{\degres}{\ensuremath{^\circ}}
\providecommand{\degree}{\ensuremath{^\circ}}
\providecommand{\vv}[1]{\overrightarrow{#1}}
\providecommand{\Coord}[1]{#1}
\providecommand{\decosix}{}
\providecommand{\starredbullet}{\ensuremath{\bullet}}
\pagestyle{empty}
\setlength{\parindent}{0pt}
\usepackage{adjustbox}\DeclareUnicodeCharacter{2236}{:}\DeclareUnicodeCharacter{2190}{\ensuremath{\leftarrow}}
\providecommand{\R}{\mathbb{R}}
\providecommand{\N}{\mathbb{N}}
\providecommand{\Z}{\mathbb{Z}}
\providecommand{\Q}{\mathbb{Q}}
\providecommand{\D}{\mathbb{D}}
\providecommand{\C}{\mathbb{C}}
\providecommand{\e}{{\mathrm{e}}}
\providecommand{\vect}[1]{\overrightarrow{#1}}
\providecommand{\vectt}[1]{\overrightarrow{\mathrm{#1}}}
\providecommand{\vvt}[1]{\overrightarrow{\mathrm{#1}}}
\providecommand{\np}[2][]{{#2~#1}}
\providecommand{\vv}[1]{{\overrightarrow{#1}}}
\renewcommand{\d}{{\mathrm{d}}}
\providecommand{\degres}{{^\circ}}
\providecommand{\degre}{{^\circ}}
\providecommand{\blue}{\color{blue}}
\providecommand{\red}{\color{red}}
\providecommand{\enskip}{\quad}
\providecommand{\llbracket}{\lbrack\!\lbrack}
\providecommand{\rrbracket}{\rbrack\!\rbrack}
\providecommand{\euro}{{\text{€}}}
\providecommand{\no}{{\text{n°}}}
\providecommand{\interval}[2]{{[#1;#2]}}
\providecommand{\intervFO}[2]{{[#1;#2[}}
\providecommand{\dotfill}{\ldots}

\providecommand{\ds}{\displaystyle}
\providecommand{\textding}[1]{#1}
\providecommand{\Oij}{\(\mathrm{O};\vec{\imath},\vec{\jmath}\)}
\providecommand{\Oijk}{\(\mathrm{O};\vec{\imath},\vec{\jmath},\vec{k}\)}
\providecommand{\Ouv}{\(\mathrm{O};\vec{u},\vec{v}\)}
\providecommand{\intervFF}[2]{\ensuremath{[#1;#2]}}
\renewcommand{\euro}{\ifmmode\text{\texteuro}\else\texteuro\fi}
\DeclareUnicodeCharacter{00B0}{\ifmmode^{\circ}\else\textdegree\fi}
\DeclareUnicodeCharacter{20AC}{\ifmmode\text{\texteuro}\else\texteuro\fi}
% Énoncés : Éducation nationale. Transcription et corrigés : APMEP, Denis Vergès et contributeurs. Découpage et classement : CoopMaths, Sébastien Lozano, Éric Elter et contributeurs. Adaptation de la présentation : Maths Explorer.
% CC BY-SA : https://coopmaths.fr/www/about/
% https://coopmaths.fr/alea/static/bac/2021/tex/bac_2021_06_metropole_6_cor.tex
\begin{document}
\noindent\textit{APMEP / CoopMaths. CC BY-SA. Présentation adaptée par Maths Explorer.}\par\medskip
\textbf{Commun à tous les candidats}

\medskip



























\textbf{Question 1} : 
On voit que pour $t = 5$, les coordonnées du point de la droite $\mathcal{D}'$ sont $(11~;~- 9~;~- 22)$ soit les coordonnées de M$_2$.

\smallskip

\hfill\textbf{Réponse b.}

\medskip

\textbf{Question 2} : Un vecteur directeur de la droite $\mathcal{D}'$ est : $\vect{u}_3\begin{pmatrix}3\\- 3\\-6\end{pmatrix}$.

\smallskip

\hfill\textbf{Réponse c.}

\medskip



\textbf{Question 3} : 






Un vecteur directeur de la droite $\mathcal{D}$ est $\vect{\text{AB}}\begin{pmatrix}-2\\2\\4\end{pmatrix}$ colinéaire au vecteur $\dfrac{1}{2}\vect{\text{AB}}\begin{pmatrix}-1\\1\\2\end{pmatrix}$

Un vecteur directeur de la droite $\mathcal{D}'$ est $\vect{u_3}\begin{pmatrix}3\\- 3\\-6\end{pmatrix}$ colinéaire au vecteur $\dfrac{1}{3}\vect{u_3}\begin{pmatrix}1\\- 1\\-2\end{pmatrix}$ ou encore colinéaire au vecteur $-\dfrac{1}{3}\vect{u_3}\begin{pmatrix}-1\\+ 1\\2\end{pmatrix}$.

Les droites $\mathcal{D}$ et $\mathcal{D}'$ ayant des vecteurs directeurs colinéaires au même vecteur sont donc parallèles.

De plus en remplaçant $t$ par $\dfrac{5}{3}$ dans l'équation paramétrique de $\mathcal{D}'$ on obtient $x = 1,~\, y = 1$ et $z = - 2$.

Les droites sont parallèles et ont un point commun : elles sont donc confondues.

\smallskip

\hfill\textbf{Réponse d.}

\medskip

\textbf{Question 4} : 





$\mathcal{P}$ a pour vecteur normal $\vect{p}\begin{pmatrix}1\\m\\-2\end{pmatrix}$.

$\mathcal{D}$ est parallèle au plan $\mathcal{P}$ si $\vect{\text{AB}}$ et $\vect{p}$ sont orthogonaux, soit :

$\vect{\text{AB}} \cdot \vect{p} = 0 \iff -2 \times 1 + 2 \times m + 4 \times (- 2) = 0 \iff - 2 + 2m - 8 = 0 \iff 2m = 10 \iff m = 5$. 

\smallskip

\hfill\textbf{Réponse c.}

\medskip

\bigskip

\end{document}
